\scanpage{383}
$\mathcal M_{0,3}^{!\sim}$ by $\widehat{\mathbb Z}$.

The calculations of \S 46 \uline{V} (page 407 \unsure{and following}) make it possible on the other hand to make explicit the pairs $(\unsure{u|\alpha},f_0=f)$ in terms of systems
\[
(16)\qquad
(\alpha,\beta,f)\quad\left|\quad
\begin{gathered}
\alpha\in\widehat{\widetilde{\Pi}}_{0,3}
 =\widehat{\Pi}_{0,3}\cdot\{1,\tau_\infty\}
 =\text{\unsure{inverse image|inverse image} of }\{1,\tau\}\\
\text{under }\widehat{\unsure{\mathcal G|\mathcal C}}_{0,3}\longrightarrow\mathcal T_{0,3}\\
\beta,f\in\widehat{\Pi}_{0,3}
\end{gathered}\right.
\]
\margin{to $(u,f)$ we associate $(\unsure{fu},f)$ by putting $\unsure{fu}\in\mathcal M_{0,3}^{!\sim}[0]$ in the form $u_{\alpha,\beta}$}
satisfying relations (not involving $f$)
\[
(17)\qquad
\alpha\rho(\alpha)^{-1}
 =\beta\sigma(\beta)^{-1}
 \underbrace{\varepsilon_1(\mu-1)}_{=\unsure{\ell_1|h}^{\nu}\ \text{if }\mu-1=2\nu}
\]
where $\mu\in\widehat{\mathbb Z}^{*}$ is defined \deleted{\unsure{without any}} \unsure{ambiguity|unambiguously} by $\alpha$ --- and even \unsure{from knowledge alone} \deleted{of the image} of $\xi=\underline\alpha\rho(\underline\alpha)^{-1}$, where $\underline\alpha$ is the image of $\alpha$ in $\mathrm{Gl}(2,\widehat{\mathbb Z})^{\pm1}$.

The case $f=1$ corresponds to \inserted{a} subgroup of $S_0\mathcal M_{0,3}^{!\sim}$, which is identified with the normalizer of $L_0$ in $\mathcal M_{0,3}^{!\sim}$, which we had denoted $\mathcal M_{0,3}^{!\sim}[0]$ (\S 46 \uline{V} $\ldots$) :
\[
\begin{gathered}
(18)\qquad
\underset{\substack{\Vert\ \mathrm{dfn}\\
\operatorname{Norm}_{\mathcal M_{0,3}^{!\sim}}(L_0)}}{\mathcal M_{0,3}^{!\sim}[0]}
\simeq
\left\{(\alpha,\beta,f)\,\middle|\,
\begin{array}{l}f=1,\ \alpha,\beta\\\text{satisfying }(17)\end{array}\right\}
\subset S_0\mathcal M_{0,3}^{!\sim} .
\end{gathered}
\]
If we denote by $\unsure{u_{\xi,\eta}|u_{\alpha,\beta}}$ \deleted{\unsure{or $u_\xi$}} \deleted{the automorphism} \unsure{or simply} $u_\xi$ the automorphism defined in terms of
\[
(19)\qquad
\xi=\alpha\rho(\alpha)^{-1},\qquad
\eta=\beta\sigma(\beta)^{-1}
\]
\margin{382}

% Printed archive footer (not autograph).
\begin{center}
382\\
\texttt{https://grothendieck.umontpellier.fr}\\
Université de Montpellier, UMR CNRS 5149, Case courrier 051, Place Eugène Bataillon\\
34095 Montpellier cedex 5 - France
\end{center}

\scanpage{384}
\mspage{493}
by
\[
(20)\qquad u(\rho)=\xi\rho,\qquad u(\sigma)=\eta\sigma
\]
(N.B. $\eta$ is defined in terms of $\xi$ thanks to (17), where $\mu=\mu(\xi)$ is determined by $\xi$), we \unsure{will have} by construction (loc. cit.)
\[
(21)\qquad u_{\alpha,\beta,f}=\operatorname{int}(f^{-1})u_{\alpha,\beta}
\]
and to the triple $(\alpha,\beta,f)$ corresponds the element $(u_{\alpha,\beta,f},f)$ of $S_0\mathcal M_{0,3}^{!\sim}$ :
\[
(22)\qquad (\alpha,\beta,f)\longmapsto (u_{\alpha,\beta},f)\in S_0\mathcal M_{0,3}^{!\sim} .
\]
Having thus made the elements of $S_0\mathcal M_{0,3}^{!\sim}$ explicit, we define a map
\[
(23)\qquad
\diagram{\begin{tikzcd}[ampersand replacement=\&,column sep=large,row sep=large]
S_0\mathcal M_{0,3}^{!\sim}\arrow[r] \&
S\mathcal M_{\rho,\sigma}(\widehat{\mathbb Z})'
=S\mathcal G_{\rho,\sigma}^{*}(\widehat{\mathbb Z})'
\times\mathrm{Sl}(2,\widehat{\mathbb Z})\unsure{'}\\
(\alpha,\beta,f)\arrow[r,mapsto] \&
\left(\underline\alpha,\underline\beta,
\mu=\mu(\underline\xi=\underline\alpha\rho(\underline\alpha)^{-1}),
\underset{\substack{\uparrow\\\inserted{\unsure{\varepsilon}}}}{\varepsilon_2(\mu)},\underline f\right)
\end{tikzcd}}
\]
where the underlined $\underline\alpha,\underline\beta,\underline f$ \unsure{\&} denote the images of $\alpha,\beta,f$ in $\mathrm{Sl}(2,\widehat{\mathbb Z})'$, and the \unsure{``accent''} $'$ in $S\mathcal M_{\rho,\sigma}(\widehat{\mathbb Z})'$ denotes the set of \deleted{pairs} \inserted{systems} $(\underline\alpha,\underline\beta,\mu,\varepsilon,\underline f)$
\deleted{\unsure{defined of elements of}}
\[
\in\mathrm{Gl}(2,\widehat{\mathbb Z})^{\prime\unsure{\pm1}}
\times\mathrm{Sl}(2,\widehat{\mathbb Z})'
\times\widehat{\mathbb Z}^{*}\times\{\pm1\}
\times\inserted{\mathrm{Sl}(2,\widehat{\mathbb Z})'}
\]
satisfying \inserted{\uline{in $\mathrm{Sl}(2,\widehat{\mathbb Z})$} (\uline{not in $\unsure{\mathrm{Sl}}(2,\widehat{\mathbb Z})'$ !}),} the loop relation \deleted{\unsure{taken|by} modulo passage to}
\[
(24)\qquad
\underline\alpha\rho(\underline\alpha)^{-1}
 =\varepsilon\underline\beta\sigma(\underline\beta)^{-1}\varepsilon_1(\mu-1)
\]
(which does not depend on the choices of \uline{signs} $\ldots$)
\margin{383}

% Printed archive footer (not autograph).
\begin{center}
383\\
\texttt{https://grothendieck.umontpellier.fr}\\
Université de Montpellier, UMR CNRS 5149, Case courrier 051, Place Eugène Bataillon\\
34095 Montpellier cedex 5 - France
\end{center}

\scanpage{385}
We have ``done what was needed'' so that the restriction of (23) to $\deleted{\unsure{S_0}}\mathcal M_{0,3}^{!\sim}[0]$ ((18), identified with the $(\alpha,\beta,f)$ for which $f=1$) is a homomorphism \inserted{into the group $S\mathcal G_{\rho,\sigma}$} --- by \unsure{defining} precisely in a clever way \inserted{ad hoc} the composition law in $S\mathcal G_{\rho,\sigma}^{*}$.
I am going to recheck it all the same to \unsure{reassure myself}, by considering $u=u_{\alpha,\beta}$, $u'=u_{\alpha',\beta'}$, $u''=uu'$ in the form $u_{\alpha'',\beta''}$, and recalling the ``cocycle relation'' which gives, up to one of the two \unsure{signs}, $(\alpha'',\beta'')$ in terms of $(\alpha',\beta'),(\alpha,\beta)$ (\S 46 \uline{II} $5^\circ$))
\[
(25)\qquad
\alpha''=u'_{\unsure{\infty|\omega}}(\alpha)\alpha',\qquad
\beta''=u'(\beta)\beta',\qquad
\inserted{\text{in addition to }\mu''=\mu\mu'}
\]
Let \unsure{then|therefore} $\underline\alpha,\underline\beta,\underline\alpha',\underline\beta'$ be \deleted{images} \inserted{\unsure{elements}} of $\mathrm{Gl}(2,\widehat{\mathbb Z})^{\pm1}$ (determined up to \uline{sign}) which lift the images of $\alpha,\alpha',\beta,\beta'$ in $\mathrm{Gl}(2,\widehat{\mathbb Z})'$, we will therefore have $(\underline\alpha,\underline\beta,\mu,\varepsilon_2(\mu))$ and $(\underline\alpha',\underline\beta',\mu',\varepsilon_2(\mu'))$ in $S\mathcal G_{\rho,\sigma}^{*}(\widehat{\mathbb Z})$, and applying $\Theta$ we find $\underline u,\underline u'\in\mathrm{Gl}(2,\widehat{\mathbb Z})$ (in fact, $\in B'_0(\widehat{\mathbb Z})$) such that
\[
(26)\qquad
\left\{\begin{aligned}
\underline u'(\rho)&=\operatorname{int}(\underline\alpha')\rho
&\underline u'(\sigma)&=\varepsilon_2(\mu')\operatorname{int}(\underline\beta')(\sigma)\\
\det\underline u'&=\mu'
\end{aligned}\right.
\]
(and relations \unsure{item|idem} for $u$)
and we will then have, for every $x\in\deleted{\unsure{\Pi}}\widehat{\unsure{\mathcal G|\mathcal C}}_{0,3}$
\[
(27)\qquad \underline{u'(x)}=\underline u'(\underline x)
\]
(cf. $\ldots$) Thus, for \deleted{some} suitable liftings $\underline\alpha'',\underline\beta^{\unsure{''|'}}$ of the images of $\alpha^{\unsure{''|'}},\beta^{\unsure{''|'}}$ in $\mathrm{Gl}(2,\widehat{\mathbb Z})^{\pm1}$, we will have
\[
\begin{gathered}
(28)\qquad
\underline\alpha''=\underline u'(\underline\alpha)\underline\alpha',\qquad
\underline\beta''=\underline u'(\underline\beta)\underline\beta',\\
\text{in addition to }\mu''=\mu\mu'\quad\text{hence}\quad
\varepsilon_2''=\varepsilon_2\varepsilon_2'.
\end{gathered}
\]
\margin{384}

% Printed archive footer (not autograph).
\begin{center}
384\\
\texttt{https://grothendieck.umontpellier.fr}\\
Université de Montpellier, UMR CNRS 5149, Case courrier 051, Place Eugène Bataillon\\
34095 Montpellier cedex 5 - France
\end{center}

\scanpage{386}
\mspage{494}
and it remains to check that the element
\[
\begin{gathered}
g''=(\underline\alpha'',\underline\beta'',\mu''=\mu\mu',\varepsilon_2''=\varepsilon_2'\varepsilon_2)
\in\deleted{\unsure{\Gamma}(\mathrm{Gl}(2,\widehat{\mathbb Z})\times\mathrm{Sl}(2,\widehat{\mathbb Z})\times}\\
\inserted{\Gamma\bigl(\widehat{\mathbb Z};\mathrm{Gl}(2)\times\mathrm{Sl}(2)\times\unsure{\mathbf E}'\times\mu_2\bigr)}
\end{gathered}
\]
is indeed the product of the elements
\[
g'=(\underline\alpha',\underline\beta',\mu',\varepsilon')\quad\text{and}\quad
g=(\underline\alpha,\underline\beta,\mu,\varepsilon).
\]
That it is OK for $\mu''$ and $\varepsilon''$ is trivial, it remains to check that in $\mathcal G_{\rho,\sigma}^{*}$, the product is indeed (as regards the components in $\alpha$ and in $\beta$) given by formulas (28) --- which \unsure{will be} valid, of course, over any base, and under the general conditions of \uline{V} (or \inserted{even} of \uline{III}, if we are interested in $\mathcal G_{\rho,\sigma}$, instead of $\mathcal G_{\rho,\sigma}^{*}$). We will have
\[
(29)\qquad
\left\{\begin{aligned}
\underline\alpha&=\underline u\,\underline a^{-1},&
\underline\beta&=\underline u\,\underline b^{-1}\\
\underline\alpha'&=\underline u'\,\underline a'^{-1},&
\underline\beta'&=\underline u'\,\underline b'^{-1}
\end{aligned}\right.
\qquad
\begin{gathered}
\underline\alpha,\underline\alpha'\in T_\rho(\widehat{\mathbb Z})\\
\underline\beta,\underline\beta'\in N'_\sigma(\widehat{\mathbb Z})
\end{gathered}
\]
and by definition of the products in $\mathcal G_{\rho,\sigma}^{*}$, the product of \deleted{\unsure{computed at}} $g'g$ will therefore have components $\underline\alpha_0''$ and $\underline\beta_0''$
\[
\underline\alpha_0''=(\underline u'\underline u)(\underline a'\underline a)^{-1},\qquad
\underline\beta_0''=(\underline u'\underline u)(\underline b'\underline b)^{-1}
\]
and it remains to prove
\[
\underline\alpha_0''=\underline\alpha'',\qquad
\underline\beta_0''=\underline\beta'',\quad\text{i.e. the relations}
\]
\[
(30)\qquad
\underbrace{\underline u'(\underline\alpha)\cdot\underline\alpha'}
 =(\underline u'\underline u)(\underline a'\underline a)^{-1}
\qquad
\underline u'(\underline\beta)\underline\beta'
 =(\underline u'\underline u)(\underline b'\underline b)^{-1}
\]
where $\deleted{\unsure{\underline u}}\ \underline\alpha,\underline\beta,\underline\alpha',\underline\beta'$ are given in terms of \deleted{$\unsure{\underline\alpha,\underline\beta,\underline u,\underline a}$} $\underline u,\underline a,\underline b,\underline u',\underline a',\underline b'$ by (29). We therefore have
\[
\begin{aligned}
\underline u'(\underline\alpha)\underline\alpha'
 &=\underbrace{\underline u'(\underline u\,\underline a^{-1})\underline u'^{-1}}\underline u'\underline a'^{-1}\\
 &=(\underline u'\underline u)(\underline a^{-1}\underline a'^{-1})
 =(\underline u'\underline u)(\underline a'\underline a)^{-1},
 \qquad\deleted{\text{\unsure{OK}}}
\end{aligned}
\]
\deleted{$\unsure{\underline u'(\underline\beta)\underline\beta'}$} and likewise for the second relation 30, which proves our assertion.
\margin{$\unsure{\checkmark|/}$}
\margin{385}

% Printed archive footer (not autograph).
\begin{center}
385\\
\texttt{https://grothendieck.umontpellier.fr}\\
Université de Montpellier, UMR CNRS 5149, Case courrier 051, Place Eugène Bataillon\\
34095 Montpellier cedex 5 - France
\end{center}